\section{Related Work}
\label{sec:related}

\subsection{Step-indexing and the topos of trees}
\label{sec:topos}

Step-indexed logical relations were introduced by Appel and McAllester
\cite{appel-mcallester2001} to give a semantic account of recursive
types: two terms are equivalent at index $n$ if no context running for
at most $n$ steps can distinguish them, and the desired equivalence is
the intersection over all $n$.  The downward-closure property they
identify is the monotonicity of step-indexed propositions in our
framework.  Ahmed \cite{ahmed2004} extended the technique to
higher-order store, the canonical setting where step-indexing is
essential rather than convenient.  Birkedal and collaborators
\cite{birkedal-et-al-2011-stepindexed} scaled it to capability calculi.
The \emph{later} modality $\triangleright$ itself was introduced by
Appel, Melli\`es, Richards, and Vouillon \cite{appel-melham-mcallester2007}
to guard the step-indexed recursion syntactically.

The categorical reformulation is due to Birkedal, Møgelberg,
Schwinghammer, and Støvring \cite{birkedal-et-al-2012}: step-indexed
propositions are the propositions of the internal logic of the
\emph{topos of trees}, the presheaf topos on $(\omega, \le)$, with
$\triangleright$ as a particular endofunctor and Löb's rule reflecting
well-foundedness of the index domain.  Subsequent work refines the
categorical setting and uses it to model richer languages
\cite{bizjak-grathwohl-clouston-mogelberg-birkedal-2016}.  A presheaf
topos on $(P, \le)$ is the topos of sheaves for the trivial
Grothendieck topology on $P$, the same data as a Kripke model of
intuitionistic logic with worlds in $P$, and hence implicitly a
forcing semantics in the model-theoretic sense.  The word
``forcing'' is not, to our knowledge, prominent in the topos-of-trees
literature. Making the terminological connection explicit is one of
the goals of this paper.

The current state of the art on the PL side is the \emph{Iris} family
of separation logics
\cite{jung-et-al-2018,krebbers-jung-bizjak-birkedal-dreyer-2017} with
extensions to higher-order ghost state, transfinite step-indices
\cite{spies2021transfinite}, and a substantial Coq mechanization.  The
internal logic of Iris is a step-indexed modal logic with $\triangleright$
and a Löb rule used pervasively as a workhorse.  Related step-indexed
verification systems include VST \cite{appel2014vst}.  Atkey and
McBride \cite{atkey-mcbride2013} use $\triangleright$ to characterize
productive coprogramming, and Clouston, Bizjak, Grathwohl, and
Birkedal \cite{clouston-et-al-2015,clouston-et-al-2017} develop
guarded $\lambda$-calculi with $\triangleright$ as a primitive type
former. Møgelberg and others \cite{mogelberg2014,bizjak-mogelberg-2018}
extend these to dependent types.

\paragraph{Why not build this on Iris?}
A natural question is whether our results could be stated as facts
inside Iris.  For the diagnostic of Section~\ref{sec:imp-vs-lam}, the
answer is no in practice: Iris's soundness theorem is monolithic and
its internal Löb rule is used whenever a recursive predicate appears,
so the distinction between recursive Hoare rules where the soundness
predicate's recursion is in the depth versus in the program is
invisible inside the Iris idiom.  Iris does the right thing because
its rule subsumes both cases, but the fact that it \emph{has} to use
the internal rule in one case and merely \emph{may} use it in the
other is not articulated in the Iris codebase or documentation that
we have been able to find.  For the counterexample of
Section~\ref{sec:lob-essential}, Iris's model builds in
well-foundedness at the level of its uPred construction, so the
question ``what fails if the step-index domain is not well-founded?''
cannot be asked inside the framework.  The minimal framework of this
paper is intended not to compete with Iris but to expose the
structural features of its soundness machinery in a setting small
enough that the two observations are stateable.

\subsection{Forcing, ramified and unramified, and the JTS translation}

Cohen introduced forcing in 1963
\cite{cohen1963-independence-i,cohen1963-independence-ii,cohen1966}
to prove the independence of the continuum hypothesis from
\textsf{ZFC}.  The original presentation is \emph{ramified forcing}:
the construction of the generic extension is stratified by ordinal
levels mimicking Gödel's constructible hierarchy, with a
``ramified language'' having a type symbol for each ordinal below
some bound and a forcing relation $p \Vdash_\alpha \varphi$ defined
level by level.  The two features that make this presentation
``ramified'', predicative stratification of quantification and
explicit level descent, are inherited from Russell and Whitehead's
ramified type theory \cite{russell-whitehead-PM} and the predicative
analysis tradition \cite{weyl1918,feferman1998}.  Shoenfield's 1971
reformulation \cite{shoenfield1971} eliminated the syntactic
ramification: forcing conditions, $P$-names, and the forcing
relation are defined by recursion on rank and formula complexity
directly.  Shoenfield's presentation is equivalent and is the
modern textbook standard
\cite{kunen2011,jech2003,scott-solovay1967}.  What was eliminated,
the explicit ramified syntactic hierarchy, is exactly what we
argue is \emph{not} eliminable in the step-indexed setting.

A parallel literature in the constructive type theory community
develops syntactic \emph{forcing translations} of type theories.
Jaber, Tabareau, and Sozeau \cite{jaber-tabareau-sozeau-2012} give a
forcing translation of CIC: for a forcing notion $P$, a syntactic
translation takes a CIC term $M : T$ to a forced term $M^P : T^P$.
For propositions, $T^P$ is a downward-closed predicate on $P$:
exactly our \texttt{iProp}.  The line was extended in
\cite{jaber-et-al-2016} and is related to work by Coquand and
collaborators on constructive forcing in type theory
\cite{coquand-jaber2010}.

The JTS translation is also available as a Coq plugin
\cite{forcing-coq-plugin}, a translation tool that develops no Hoare
logic, wp predicate, or operational-semantics soundness proof.
Appendix~\ref{sec:bridge} compares it in detail.  What the JTS line
provides and we do not is a full translation of CIC including
dependent types.  The propositional fragment of that translation is
formalized against our \texttt{iProp} in Appendix~\ref{sec:bridge}
(the factorization theorem \texttt{jts\_is\_interp}); the term-level
fragment is the natural future direction.

\subsection{The modal logic of provability}

The Gödel--Löb provability logic \textsf{GL} is the modal logic
characterized by the axiom $\Box(\Box \varphi \to \varphi) \to
\Box \varphi$.  Solovay's arithmetical completeness theorem
\cite{solovay1976} establishes that \textsf{GL} is sound and complete
for arithmetical interpretations of $\Box$ as ``provable in PA.''  On
the model-theoretic side, \textsf{GL} is sound and complete for the
class of transitive Kripke frames whose accessibility relation is
\emph{converse well-founded}. See Smoryński \cite{smorynski1985} and
Boolos \cite{boolos1993}.

The relevance to our work is direct.  The $\triangleright$ modality of
step-indexed logic is formally a \textsf{GL}-style $\Box$ once the
strict refinement order $\prec$ on conditions is identified with the
converse of the \textsf{GL} accessibility relation.  Löb's rule of
step-indexed logic is the Löb axiom of \textsf{GL}.  The requirement
that the forcing notion be well-founded is the dual of \textsf{GL}'s
requirement that the accessibility relation be converse well-founded.
Our Section~\ref{sec:lob-essential} result is a special case of the
classical observation that \textsf{GL} is not sound on frames where
that condition fails.  In modal logic this is folklore.  Brought into
the step-indexed Hoare logic setting and mechanized, as far as we can
determine, it has not previously appeared.

\subsection{What this paper contributes}

Each ingredient is individually known to specialists in some adjacent
community: topos-of-trees to guarded domain theory, \textsf{GL}
model theory to modal logicians, ramified forcing to set theorists,
the JTS translation to type theorists.  What does not appear as a
single mechanized account is a step-indexed Hoare logic that
articulates the recursion-locus diagnostic of
Section~\ref{sec:imp-vs-lam} and exhibits the well-foundedness
counterexample of Section~\ref{sec:lob-essential}, together with the
forcing-style re-presentation that makes both stateable.  We omit a
survey of classical Hoare and separation logic, referring the reader
to standard sources
\cite{hoare1969,ohearn-reynolds-yang-2001,brookes2007,cook1978}.
